% Blueprint note: Fixed upstream source versions: \cite{KIPSourcebbac1fc05843}. Interface scope notes are retained as native hgraph node comments; historical migration records are local archives.

% Curated mathematical exposition; compiled dependencies and status are supplied by the publisher.

\section{Common representatives and page calculus}


\begin{definition}[Common cycle, boundary and permanent labels]
\label{sem:ProofTools:common-labels}
\level{coarse}
For a classical spectrum $X$, a common label in bidegree $(s,t)$ is an element of $E_2^{s,t}(X)$. For $c\geq1$, it is a level-$c$ cycle, respectively boundary, if it lies in $Z_c^{s,t}(X)$, respectively $B_c^{s,t}(X)$. A permanent label is an element of $Z_\infty^{s,t}(X)$; the zero label is included.
\lean{KIP126.ProofTools.Ambient,KIP126.ProofTools.IsCycle,KIP126.ProofTools.IsBoundary,KIP126.ProofTools.IsPermanent}
\end{definition}

\begin{definition}[Congruence and differential equations]
\label{sem:ProofTools:congruence-and-differential}
\level{coarse}
For $c\geq1$, labels $x,y\in E_2^{s,t}(X)$ are congruent at level $c$ if $x-y\in B_c^{s,t}(X)$. For $r\geq2$, the equation $d_r(x)=y$ means that $x$ and $y$ represent classes on page $r$ and their classes satisfy the Adams differential equation, with source bidegree $(s,t)$ and target bidegree $(s+r,t+r-1)$. The target may be zero.
\lean{KIP126.ProofTools.Congruent,KIP126.ProofTools.DifferentialAt}
\end{definition}

\begin{theorem}[Cycles and representatives on later pages]
\label{sem:ProofTools:cycle-reaches}
\level{coarse}
For $r\geq2$, a common label lies in $Z_{r-1}$ if and only if it reaches page $r$. If $x$ represents $x_r$ on page $r$ and $y-x\in B_{r-1}$, then $y$ represents $x_r$. If $d_r(x)=y$ and $y'-y\in B_{r-1}$, then $d_r(x)=y'$.
\lean{KIP126.ProofTools.isCycle_iff_reaches,KIP126.ProofTools.representsOnPage_of_congruent,KIP126.ProofTools.differentialAt_of_congruent}
\end{theorem}

\begin{theorem}[Boundary and cycle calculus]
\label{sem:ProofTools:boundary-calculus}
\level{coarse}
For $c\geq1$, $B_c\subseteq Z_c$, and $B_1=0$. If a common label represents zero on page $r\geq2$, it lies in $B_{r-1}$. Projection of a common cycle to the initial page recovers its $E_2$-label.
\lean{KIP126.ProofTools.Internal.Mahowald.boundaries_le_cycles,KIP126.ProofTools.Internal.Mahowald.boundaries_one,KIP126.ProofTools.Internal.Mahowald.represents_zero_mem_boundary,KIP126.ProofTools.Internal.Mahowald.initialProjection_eq_cycleLabel}
\end{theorem}

\begin{theorem}[Composition and infinite representatives]
\label{sem:ProofTools:page-naturality}
\level{coarse}
\uses{sem:AdamsSynthetic:families,sem:SpectralSequence:permanent}
For composable synthetic maps $f$ and $g$, their induced page maps satisfy $(g\circ f)_r=g_r\circ f_r$. If $x$ represents an infinite-page class $e$, then $f_2(x)$ represents $f_\infty(e)$.
\lean{KIP126.ProofTools.Internal.family_page_comp,KIP126.ProofTools.Internal.family_infinity_representative_map}
\end{theorem}

\section{Finite filtration arguments}


\begin{theorem}[Finite correction of a residual]
\label{sem:ProofTools:finite-correction}
\level{coarse}
Let $A,A',B$ be abelian groups, $N\in\mathbb N$, and $\rho:A'\to A$ and $g:A\to B$ homomorphisms, let $E_k\leq A$, $L_k\leq A'$ with $L$ decreasing, and let $T\leq B$. Suppose that for each $s\leq k<s+N$, a residual $\delta\in E_k$ can either be corrected by some $\varepsilon\in L_k$ so that $\delta-\rho\varepsilon\in E_{k+1}$, or already satisfies $g\delta\in T$. Suppose also that $g(E_{s+N})\subseteq T$. If $\alpha_0-\rho\alpha\in E_s$, there is $\alpha'$ with
\[\alpha'-\alpha\in L_s,\qquad g(\alpha_0-\rho\alpha')\in T.\]

% Blueprint note: Finite correction concerns a finite interval; it does not assert a completion or a compatible infinite lift.
\lean{KIP126.ProofTools.finite_filtration_correction}
\end{theorem}

\begin{theorem}[Reflection along a finite filtration interval]
\label{sem:ProofTools:finite-reflection}
\level{coarse}
Let $B',B$ be abelian groups, $N\in\mathbb N$, and $\rho:B'\to B$ a homomorphism. Let $E$ be a decreasing filtration on $B$, and let $L_k\leq B'$. Suppose that for $s\leq k<s+N$, $b\in L_k$ and $\rho b\in E_{k+1}$ imply $b\in L_{k+1}$. Then $b\in L_s$ and $\rho b\in E_{s+N}$ imply $b\in L_{s+N}$.
% Blueprint note: Ambient-map injectivity is not required.
\lean{KIP126.ProofTools.finite_filtration_reflection}
\end{theorem}

\begin{theorem}[The first nonzero grade in a finite interval]
\label{sem:ProofTools:leading-grade}
\level{coarse}
Let $E_k$ be subgroups of an abelian group and let $N\in\mathbb N$. If $b\in E_s$ but $b\notin E_{s+N}$, then there exists $s\leq k<s+N$ with
\[b\in E_k\quad\text{and}\quad b\notin E_{k+1}.\]

% Blueprint note: No separation or completeness assumption is required.
\lean{KIP126.ProofTools.first_nonzero_filtration_grade}
\end{theorem}

\section{Convergence, filtration and detection}


\begin{definition}[The infinite-page graded identification]
\label{sem:ProofTools:infinity-graded}
\level{coarse}
\uses{sem:AdamsSynthetic:selected-convergence,sem:SpectralSequence:filtration}
For a convergent synthetic Adams tower on $X$, let
\[E_\infty^{s,m+s,w}(X)\cong\operatorname{gr}^s\pi_{m,w}X\]
be its convergence isomorphism.
\lean{KIP126.ProofTools.Internal.infinityGradedIso}
\end{definition}

\begin{theorem}[Filtration representatives and preservation]
\label{sem:ProofTools:filtration-representatives}
\level{coarse}
\uses{sem:AdamsSynthetic:selected-convergence,sem:SpectralSequence:filtration}
Every element of $F^s\pi_{m,w}X$ has a representative in the filtration subobject $F^sX$. If $g:X\to Y$ is a map of convergent towers, then
\[g_*\bigl(F^s\pi_{m,w}X\bigr)\subseteq F^s\pi_{m,w}Y.\]
\lean{KIP126.ProofTools.Internal.filtration_representative_exists,KIP126.ProofTools.Internal.map_preserves_filtration}
\end{theorem}

\begin{theorem}[Representatives for a tower extension equation]
\label{sem:ProofTools:extension-representatives}
\level{coarse}
\uses{sem:ProofTools:infinity-graded,sem:ProofTools:filtration-representatives,paper:2.3}
Let $g:X\to Y$ be a map of convergent towers. Fix $(m,w)$ and classes $x\in E_\infty^{s,m+s,w}(X)$ and $y\in E_\infty^{s+n,m+s+n,w}(Y)$. An extension equation of length $n$ from $x$ to $y$ holds if and only if $n\geq0$ and there are representatives
\[a\in F^s\pi_{m,w}X,\qquad b\in F^{s+n}\pi_{m,w}Y,\qquad g_*(a)=b\]
whose associated-graded classes correspond to $x$ and $y$ under the convergence isomorphisms.
\lean{KIP126.ProofTools.Internal.tower_extension_iff}
\end{theorem}

\begin{theorem}[Naturality and the next-filtration ambiguity]
\label{sem:ProofTools:grade-naturality}
\level{coarse}
\uses{sem:ProofTools:infinity-graded,sem:ProofTools:filtration-representatives}
Let $g:X\to Y$ be a map of convergent towers, $a\in F^s\pi_{m,w}X$ and $b\in F^s\pi_{m,w}Y$. If $g_*(a)=b$, their associated-graded classes satisfy $g_\infty([a])=[b]$. Conversely, this graded equality implies
\[b-g_*(a)\in F^{s+1}\pi_{m,w}Y.\]
If $g_\infty$ is injective in this degree and $g_*(a)\in F^{s+1}\pi_{m,w}Y$, then $a\in F^{s+1}\pi_{m,w}X$.
\lean{KIP126.ProofTools.Internal.grade_natural,KIP126.ProofTools.Internal.grade_difference,KIP126.ProofTools.Internal.next_filtration_of_infinity_injective}
\end{theorem}

\begin{theorem}[Transport and reflection of detection]
\label{sem:ProofTools:detection-transport}
\level{coarse}
\uses{sem:ProofTools:page-naturality,sem:ProofTools:grade-naturality,sem:SpectralSequence:detection}
Let $g:X\to Y$ be a map of convergent towers. If $x$ detects $\alpha$ on $X$, then $g_2(x)$ detects $g_*(\alpha)$ on $Y$. Conversely, suppose that $g_\infty$ is injective in filtration $s$, that $x$ represents an infinite-page class, and that $\alpha\in F^s\pi_{m,w}X$. If $g_2(x)$ detects $g_*(\alpha)$, then $x$ detects $\alpha$.
\lean{KIP126.ProofTools.Internal.detects_map,KIP126.ProofTools.Internal.detects_reflect}
\end{theorem}

\begin{theorem}[A vanishing strip identifies filtration bounds]
\label{sem:ProofTools:filtration-gap}
\level{coarse}
\uses{sem:ProofTools:infinity-graded}
Fix a convergent synthetic Adams tower on $X$, and let $s\leq t$. If every infinite-page component $(j,m+j,w)$ vanishes for $s\leq j<t$, then for every $\alpha\in\pi_{m,w}X$,
\[\operatorname{fil}(\alpha)\geq s\quad\Longleftrightarrow\quad\operatorname{fil}(\alpha)\geq t.\]

% Blueprint note: The convergence is that of the specified synthetic tower.
\lean{KIP126.ProofTools.Internal.Mahowald.filtration_iff_of_infinity_zero}
\end{theorem}

\section{Finite lambda quotients and sphere labels}


\begin{theorem}[The finite quotient strip]
\label{sem:ProofTools:finite-strip}
\level{coarse}
\uses{sem:Literature:infinite-finite-window,sem:AdamsSynthetic:finite-quotient-separation,sem:ProofTools:filtration-gap}
For $q>0$, the infinite page of $\nu X/\lambda^q$ vanishes outside $0\leq t-w<q$. Every $\alpha\in\pi_{t-s,t}(\nu X/\lambda^q)$ has filtration at least $s$. If $q\leq m+s-w$, an element of $\pi_{m,w}(\nu X/\lambda^q)$ in filtration at least $s$ is zero.
\lean{KIP126.ProofTools.Internal.Mahowald.finite_infinity_zero_outside,KIP126.ProofTools.Internal.Mahowald.finite_bottom_filtration,KIP126.ProofTools.Internal.Mahowald.finite_high_filtration_zero}
\end{theorem}

\begin{theorem}[Injectivity of quotient reduction in the retained strip]
\label{sem:ProofTools:rho-injectivity}
\level{coarse}
\uses{sem:Literature:infinite-restriction,sem:Literature:infinite-finite-window}
For $0<i\leq j$ and $t-w<i$, the quotient reduction $\rho:\nu X/\lambda^j\to\nu X/\lambda^i$ induces an injective map on the infinite-page component in tridegree $(s,t,w)$.
% Blueprint note: The range bound refers to the smaller quotient.
\lean{KIP126.ProofTools.Internal.rho_infty_injective}
\end{theorem}

\begin{theorem}[Finite target labels]
\label{sem:ProofTools:finite-target-label}
\level{coarse}
\uses{sem:Literature:infinite-finite-window,sem:AdamsSynthetic:shift-representatives}
For $q>0$ and $0\leq k<q$, let $y\in Z_{q-k}^{s,t}(X)$. Its finite target label at weight $t-k$ represents $e\in E_\infty^{s,t,t-k}(\nu X/\lambda^q)$ if and only if the finite-window comparison carries $e$ to the class of $y$ in
\[Z_{q-k}^{s,t}(X)/B_{1+k}^{s,t}(X).\]
\lean{KIP126.ProofTools.Internal.Mahowald.finite_target_label_iff}
\end{theorem}

\begin{theorem}[Sphere labels and finite lambda detection]
\label{sem:ProofTools:sphere-label-detection}
\level{coarse}
\uses{sem:AdamsSynthetic:e2-comparison,sem:AdamsSynthetic:first-quotient-shifts,sem:Literature:infinite-lambda,sem:ProofTools:finite-target-label,sem:ProofTools:detection-transport}
The first-quotient comparison carries the first label of $y\in E_2^{s,t}(S^0)$ to $y$. Let $\alpha\in\pi_{t-s,t}(\nu S^0)$ have first-quotient image equal to this label, using the canonical zero-shift identification. For $q>0$ and $k\geq0$, the finite target label of $y$ in tridegree $(s,t,t-k)$ detects the image of $\lambda^k\alpha$ in $\nu S^0/\lambda^q$.
\lean{KIP126.ProofTools.Internal.Mahowald.sphere_first_quotient_label,KIP126.ProofTools.Internal.Mahowald.sphere_finite_lambda_detects}
\end{theorem}

\begin{theorem}[Suspension and composition coherence]
\label{sem:ProofTools:shift-coherence}
\level{coarse}
\uses{sem:AdamsSynthetic:shift-coherence,sem:AdamsSynthetic:ess-shifts}
Assume that suspension and bigraded shifts are coherent with the homotopy functors. The canonical isomorphism
\[\Sigma^{1,-q}Q\cong\Sigma(\Sigma^{0,-q}Q)\]
identifies the induced suspension and bigraded-shift maps. The inverse and composition isomorphisms for shifts induce the corresponding inverse and composition maps on homotopy groups.
\lean{KIP126.ProofTools.Internal.Mahowald.biSuspensionIso,KIP126.ProofTools.Internal.Mahowald.susp_invariance_add,KIP126.ProofTools.Internal.Mahowald.inverse_associativity,KIP126.ProofTools.Internal.Mahowald.shift_hom_associativity}
\end{theorem}

\section{Generalized Leibniz and Mahowald laws}


\begin{definition}[The generalized Leibniz law]
\label{sem:ProofTools:leibniz-law}
\level{coarse}
Let $f:X\to Y$ and write $e=e(f)$. The generalized Leibniz law is the following implication. Assume $2\leq p\leq r$, $e\leq m\leq p-2+e$, $e\leq l$, and
\[x\in Z_{r-1}^{s,t}(X),\qquad y\in Z_{r-1-m+e}^{s+m,t+m}(Y),\]
\[x_\infty\in Z_\infty^{s+r,t+r-1}(X),\qquad y_\infty\in Z_\infty^{s+r+l,t+r+l-1}(Y).\]
Suppose that
\[d_r(x)=x_\infty,\qquad d_m^{f,E_p}(x)=y,\qquad d_l^{f,E_\infty}(x_\infty)=y_\infty.\]
Assume either that the first differential has no crossing on page $p$ or that the finite extension has no crossing, and assume that the infinite extension has no crossing. Then
\[d_{r+l-m}(y)=y_\infty.\]
\lean{KIP126.ProofTools.GeneralizedLeibnizLaw}
\end{definition}

\begin{definition}[Actual essential no-crossing]
\label{sem:ProofTools:mh-actual-no-crossing}
\level{coarse}
\uses{sem:AdamsSynthetic:selected-convergence,sem:SpectralSequence:filtration}
Fix a map $F:U\to V$ of convergent synthetic Adams towers and one stem and weight $(p,w)$. For a source filtration $s$ and length $a$, actual no-crossing means
\[
\text{there is no essential }d_{v-u}^{F}(\xi)=\eta
\quad\text{with }s<u\leq v\leq s+a
\]
in this same map and bidegree. Essentiality is taken in the actual extension spectral sequence, including all shorter extension images. The diagonal $u=v$ is included.
\lean{KIP126.ProofTools.Internal.Mahowald.ActualNoCrossing}
\end{definition}

\begin{theorem}[A common source from actual no-crossing]
\label{sem:ProofTools:mh-common-source}
\level{coarse}
\uses{sem:ProofTools:mh-actual-no-crossing,sem:ProofTools:filtration-representatives,sem:ProofTools:finite-correction}
Let $F:U\to V$ and $G:U\to W$ be maps of convergent towers with the same source convergence. Fix the same stem and weight, and suppose
\[
d_a^F(\xi)=\eta,\qquad d_b^G(\xi)=\zeta.
\]
If either equation has actual no-crossing in its full higher-source triangle, then there are actual representatives $u,v,w$ of $\xi,\eta,\zeta$, respectively, with
\[F(u)=v,\qquad G(u)=w.\]
The representatives lie in filtrations $s,s+a,s+b$. Zero labels, inessential input equations, and length zero are allowed.
\lean{KIP126.ProofTools.Internal.Mahowald.actual_common_source_representative}
\end{theorem}

\begin{definition}[The strong finite-layer Mahowald assertion]
\label{sem:ProofTools:mahowald-law}
\level{coarse}
\uses{sem:ProofTools:mh-actual-no-crossing,paper:3.20,sem:ProofTools:finite-target-label}
Fix one classical triangle $X\xrightarrow fY\xrightarrow gZ\xrightarrow h\Sigma X$ and its compatible normalized synthetic triangle. Write $Q_DW=\nu W/\lambda^D$, retaining the fixed normalization and suspension comparisons in every map. Put
\[
n_1=n-e_f\geq1,\quad m_1=m-e_g\geq0,\quad k=l-e_h\geq0,\quad r=n+m+l,
\qquad e_f+e_g+e_h=1.
\]
For $i,j\geq0$ with $j<n_1$, set
\[
A=n_1+k+i-j,\quad B=m_1+j+1,\quad N=A+B=r+i,
\quad K=i+k,\quad J=A+j,\quad Q=N-K=r-k.
\]
Thus $0\leq i,K<A$, $0\leq j<B$, and $0\leq J<N$. Take classical labels in the prescribed bidegrees
\[
\bar x\in Z_{r-1}^{s,t}(Z),\quad \bar y\in Z_\infty^{s+r,t+r-1}(Z),\quad
x\in Z_{n_1-j}^{s+l,t+l-1}(X),\quad y\in Z_{m_1+1}^{s+l+n,t+l+n-1}(Y),
\]
with $d_r\bar x=\bar y$. Write the normalized vertices as $U=\Sigma^{0,e_f}\nu X$, $V=\nu Y$, and $W=\Sigma^{0,-e_g}\nu Z$. Choose actual suspension--quotient isomorphisms
\[
\varepsilon_X^D:Q_D(\Sigma U)\cong\Sigma Q_DU\quad(D=A,N),\qquad
\varepsilon_Y^N:Q_N(\Sigma V)\cong\Sigma Q_NV.
\]
Each must commute with its quotient inclusion. Compatibility with reduction $Q_N\to Q_A$ on bigraded sphere representatives follows internally from these inclusion squares and the finite no-return comparison; it is not an extra hypothesis. The assertion below is quantified over these choices. The map $h_A$ uses $\varepsilon_X^A$, and the output uses $\varepsilon_X^N,\varepsilon_Y^N$; a comparison preserving only objects is insufficient. Finite source labels are transported by these actual shifts and quotients, using the intrinsic finite source label rather than an independently chosen positive-shift $E_2$ identification.

Their actual finite infinite-page classes must support the following three equations in the indicated fixed maps:
\[
d_l^{h_A}(\lambda^i\bar x)=\lambda^Kx,\qquad
d_r^{\delta_{A,B}}(\lambda^i\bar x)=\lambda^{B-1}\bar y,\qquad
d_m^{g_B}(\lambda^jy)=\lambda^{B-1}\bar y.
\]
The first two equations use the same infinite-page source class. The last two use the same top-layer target class. The $h$ target is related to the specified $x$ by the fixed actual suspension comparison. Assume actual no-crossing on either $h_A$ or $\delta_{A,B}$: the entire triangle $s<u\leq v\leq s+l$ or $s<u\leq v\leq s+r$, respectively. Both are at the same actual stem and weight. In unshifted $Z$ coordinates this is $(t-s,t-i)$; in the normalized third vertex it is $(t-s,t-i-e_g)$.

The strong finite-layer Mahowald assertion is that there exist the same actual representatives $a,b,c$, with
\[
a\text{ detecting }\lambda^i\bar x,\quad h_A(a)\text{ detecting }\lambda^Kx,
\quad \delta_{A,B}(a)=g_B(b)\text{ detecting }\lambda^{B-1}\bar y,
\quad b\text{ detecting }\lambda^jy,
\]
\[
\rho(c)=h_A(a),\qquad f_N(c)=-\lambda^Ab,
\]
such that $c$ detects $\lambda^Kx$ and its actual $f_N$ image detects $\lambda^Jy$. Equivalently for the resulting page relation,
\[
x\in Z_Q,\qquad d_n^{\hat f_N}(\lambda^Kx)=[\lambda^Jy].
\]
The source label lies in $Z_Q/B_{K+1}$; the target lies in $Z_{m_1+1}/B_{J+1}$ modulo shorter actual $\hat f_N$ extension images. The actual representatives are retained in addition to this quotient relation. At $i=j=0$ the output is in the strong quotient $Q_r$:
\[
x\in Z_{n+m+e_h},\qquad d_n^{\hat f_r}(\lambda^kx)=[\lambda^{r'-1}y],\quad r'=n_1+k+1.
\]
A zero target label means a higher filtration, not a zero homotopy image. No power of $\lambda$ is cancelled.
% Blueprint note: This entry binds the definition of the full assertion. The complete theorem is bound separately at Theorem 6.12. Its proof preserves the same representatives, signed images, exact labels and terminal conclusions, and consumes the proved fixed-model Literature row internally.
% Blueprint note: The authorized minimal Literature theorem now constructs the actual signed distinguished quotient row and an inclusion-compatible suspension comparison. The full Mahowald theorem has no additional horizontal-row premise. Existing upstream construction and comparison obligations remain registered.
\lean{KIP126.ProofTools.Standard.FiniteMahowaldAtLevelStatement}
\end{definition}

\begin{theorem}[The target-preserving lifting obstruction]
\label{sem:ProofTools:mh-target-preserving-lift}
\level{coarse}
\uses{sem:ProofTools:filtration-representatives,sem:AdamsSynthetic:lambda-quotients}
Fix $\varphi_N:Q_NU\to Q_NV$ and one stem and weight. Let $p:\pi U\to\pi Q_NU$ be quotient inclusion and let $\partial$ be its actual cofiber connecting map. Choose subgroups $F\leq\pi U$, $S_+\leq S\leq\pi Q_NU$, and $T_+\leq\pi Q_NV$, with source-filtration reflection $p(\alpha)\in S\Rightarrow\alpha\in F$. For $c_0\in S$, put
\[V_0=S_+\cap\varphi_N^{-1}(T_+).\]
Then
\[
\partial c_0\in\partial V_0
\quad\Longleftrightarrow\quad
\exists\alpha\in F:\ p(\alpha)-c_0\in S_+,\quad
\varphi_Np(\alpha)-\varphi_Nc_0\in T_+.
\]
In the MH application, $S_+$ and $T_+$ are the filtrations immediately above the prescribed source and target. The original $c_0$ need not lift; some member of its target-preserving affine family must lift.
\lean{KIP126.ProofTools.lambda_quotient_target_preserving_lift_iff}
\end{theorem}

\begin{theorem}[A specified finite target from a full lift is permanent]
\label{sem:ProofTools:mh-lifted-target-permanent}
\level{coarse}
\uses{sem:ProofTools:finite-target-label,sem:ProofTools:filtration-representatives,sem:Literature:infinite-finite-window}
Let $N>0$ and $0\leq t-w<N$. Fix $y\in E_2^{s,t}(X)$ and an actual full class $\beta\in\pi_{t-s,w}\nu X$. If the projection of $\beta$ to $Q_NX$ detects the specified finite label $\lambda^{t-w}y$, then the original label $y$ belongs to $Z_\infty^{s,t}(X)$. Filtration reflection in the visible window is proved internally. No permanence assumption on $y$ is made; equality only modulo uncontrolled shorter extension images would not suffice.
\lean{KIP126.ProofTools.finite_target_from_full_lift_is_permanent}
\end{theorem}

\begin{definition}[The same-quotient nonliftable obstruction]
\label{sem:ProofTools:nonliftable-obstruction}
\level{coarse}
Let $r_{\rm early}\leq r_{\rm late}$. A nonliftable obstruction is an extension of length $n-a-b$ for the normalized map in the $(r_{\rm early}-1)$-st quotient, in stem $t-s$ and weight $t+e(f)$, where
\[a,b\in\mathbb N,\qquad 0<a\leq r_{\rm early}-2,\qquad a+b\leq n.\]
Its source has filtration $s+a$ and lies outside the image of quotient reduction from the $(r_{\rm late}-1)$-st quotient on $E_\infty$; its target has filtration $s+n-b$ and is nonzero. The extension may be inessential.
% Blueprint note: The nonliftable obstruction is stronger than an essential classical crossing.
\lean{KIP126.ProofTools.FiniteExtensionNonliftableCrossing}
\end{definition}

\begin{definition}[The finite stretching sufficient condition]
\label{sem:ProofTools:finite-stretching-law}
\level{coarse}
Suppose $2\leq r_{\rm early}\leq r_{\rm late}$ and
\[e(f)\leq n\leq r_{\rm early}-2+e(f).\]
The finite stretching law is the following implication: if
\[x\in Z_{r_{\rm late}-1}^{s,t}(X),\qquad y\in Z_{r_{\rm late}-1-n+e(f)}^{s+n,t+n}(Y),\qquad d_n^{f,E_{r_{\rm early}}}(x)=y,\]
and there is no nonliftable obstruction as above, then
\[d_n^{f,E_{r_{\rm late}}}(x)=y.\]
% Blueprint note: Finite stretching does not supply an untruncated extension or a compatible infinite tower.
\lean{KIP126.ProofTools.FinitePageExtensionStretchingLaw}
\end{definition}

\begin{theorem}[The zero-layer nu-cofiber Mahowald instance]
\label{sem:ProofTools:nu-zero-layer}
\level{coarse}
\uses{sem:Literature:nu-cofiber,sem:Literature:connected-homology,sem:ProofTools:boundary-calculus,sem:ProofTools:finite-target-label,sem:ProofTools:finite-strip,sem:ProofTools:rho-injectivity,sem:ProofTools:shift-coherence,sem:ProofTools:extension-representatives,sem:ProofTools:detection-transport,paper:5.4}
In the fixed cofiber triangle for $\nu$, take labels $x\in E_2^{8,133}(S^3)$, $y\in E_2^{11,136}(S^0)$, $\bar x\in E_2^{8,134}(C\nu)$ and permanent $\bar y\in E_2^{11,136}(C\nu)$. Require $x\in Z_2$, $y\in Z_1$, $\bar x\in Z_2$, the actual suspension label for $x$, the length-zero connecting extension on page $3$, $d_3(\bar x)=\bar y$, its specified no-crossing alternative, and the length-zero bottom-cell extension on page $2$. Then $x\in Z_3$ and there is an actual $\nu$-extension of length $3$ on page $4$ from $x$ to some $y'$ with $y'-y\in B_3$.
% Blueprint note: This specialized result has its own proof and does not require the general Mahowald theorem as an extra hypothesis.
\lean{KIP126.ProofTools.Internal.Mahowald.main_nu_mahowald_zero_layer}
\end{theorem}

\section{Synthetic boundary maps and their labels}
\begin{definition}[Synthetic boundary maps and weighted labels]
\label{sem:ProofTools:synthetic-boundary-labels}
\level{coarse}
For $n>0$, the absolute connecting map of the quotient triangle is
\[\delta_n:\nu X/\lambda^n\longrightarrow\Sigma^{1,-n}\nu X.\]
For $0<n<m$, its finite version has target $\Sigma^{1,-n}(\nu X/\lambda^{m-n})$. Using the suspension isomorphism of the quotient triangle, the source label of $x$ is $\lambda^a x$ in the $n$-th quotient. For a classical differential of length $r=k+1$ and $n\leq a+k$, its target label is the shifted $\lambda^{a+k-n}y$; the finite label is its image under the target quotient projection. A crossing of the absolute boundary on page $p$ is an essential extension in that same boundary ESS, with source filtration $u$ and target filtration $v$ satisfying $s<u\leq v\leq s+r$.
\lean{KIP126.ProofTools.Paper.syntheticAbsoluteBoundary,KIP126.ProofTools.Paper.syntheticFiniteBoundary,KIP126.ProofTools.Paper.syntheticDeltaSourceLabel,KIP126.ProofTools.Paper.syntheticBoundaryLabel,KIP126.ProofTools.Paper.syntheticDeltaNuLabel,KIP126.ProofTools.Paper.syntheticDeltaFiniteLabel,KIP126.ProofTools.Paper.DeltaHasCrossing}
\end{definition}

\section{Representative criteria and changing the page}

\begin{theorem}[Descent of finite-page extensions in the valid length range]
\label{sem:ProofTools:finite-extension-descent}
\level{coarse}
(Adapted from \cite{LWX126}, Prop. 6.20, first assertion, with the page--length domain of Def. 5.4.)\par{}

In the fixed model, let $X,Y$ be classical objects and $f:X\to Y$ a map, and put $e=e(f)$. Take finite pages $r_{\rm early},r_{\rm late}\in\mathbb N$, integers $n,s,t$, and common labels
\[
x\in E_2^{s,t}(X),\qquad y\in E_2^{s+n,t+n}(Y).
\]
Suppose that
\[
2\leq r_{\rm early}\leq r_{\rm late},\qquad
n\leq r_{\rm early}-2+e,
\]
and that there is an extension, in the sense of Definition 5.4,
\[
d_n^{f,E_{r_{\rm late}}}(x)=y.
\]
Then the same map and the same specified labels admit the earlier-page extension
\[
d_n^{f,E_{r_{\rm early}}}(x)=y.
\]
The late extension includes the lower length bound $e\leq n$ and the cycle conditions
\[
x\in Z_{r_{\rm late}-1}^{s,t}(X),\qquad
y\in Z_{r_{\rm late}-1-n+e}^{s+n,t+n}(Y).
\]
The conclusion includes the corresponding earlier-page cycle conditions. The extension need not be essential; no crossing hypothesis is imposed.
\lean{KIP126.ProofTools.Paper.finiteExtension_descends}
\uses{paper:5.4}
% Blueprint note: The explicit earlier-page upper length bound is indispensable. The proof uses the actual quotient reductions on source and target and naturality of their extension equation. This is the corrected descent statement, not the unrestricted first assertion printed in Proposition 6.20. Its proof has no local hole; registered upstream construction and proof obligations remain visible through the compiled dependency closure.
\end{theorem}


\begin{definition}[Higher-source extension ranges]
\label{sem:ProofTools:crossing-ranges}
\level{coarse}
Let $f$ be a map of convergent towers or filtered abelian groups. An extension from source filtration $s$ to target filtration $s+n$ has no higher-source crossing over $[u,s+n]$ if every extension with source filtration greater than $s$ and target filtration in $[u,s+n]$ has zero target label. All extensions, including inessential ones, occur in this quantification.
\lean{KIP126.ProofTools.Paper.ExtensionNoCrossingRange,KIP126.ProofTools.Paper.FilteredNoCrossingRange,KIP126.ProofTools.Paper.RepresentsInfinityClass}
\end{definition}

\begin{theorem}[The full representative criterion]
\label{sem:ProofTools:representative-criterion}
\level{coarse}
(From \cite{LWX126}, Prop. 2.10.)\par{}

Fix a map of convergent towers and an extension equation of length $n$ from $x$ to $y$ in filtrations $s$ and $s+n$, and a threshold $u\leq s+n$. Absence of higher-source extensions whose target lies between $u$ and $s+n$ is equivalent to the following condition: every representative $a$ of $x$ whose image has filtration at least $u$ has image representing the specified $y$ in filtration $s+n$. For threshold $u=s$, the image-filtration premise is automatic. If $y=0$, the conclusion says that every such image has filtration at least $s+n+1$. A map increasing every filtration by at least $n$ satisfies this criterion. The same assertions hold for a homomorphism of abelian groups with decreasing filtrations, using representatives of the given associated-graded source class.
\lean{KIP126.ProofTools.Paper.extension_noCrossingRange_iff,KIP126.ProofTools.Paper.extension_noCrossing_iff,KIP126.ProofTools.Paper.extension_zero_noCrossing_iff,KIP126.ProofTools.Paper.extension_noCrossing_of_filtration_bound,KIP126.ProofTools.Paper.filtered_extension_noCrossingRange_iff,KIP126.ProofTools.Paper.filtered_extension_noCrossing_iff,KIP126.ProofTools.Paper.filtered_zero_extension_noCrossing_iff,KIP126.ProofTools.Paper.filtered_extension_noCrossing_of_filtration_bound}
% Blueprint note: Detailed source attribution: (From \cite{LWX126}, Proposition 2.10; with the stated threshold formulation.)
\end{theorem}

\begin{definition}[Finite correction labels]
\label{sem:ProofTools:correction-label}
\level{coarse}
For $f:X\to Y$ and $x'\in E_2^{s+a,t+a}(X)$, the finite correction label is the image of $\lambda^a x'$ under the shift comparison and quotient isomorphism into the normalized finite source quotient, in synthetic weight $t+e(f)$.
\lean{KIP126.ProofTools.Paper.finiteExtensionCorrectionLabel}
\end{definition}

\begin{theorem}[A nonliftable short extension or a late crossing]
\label{sem:ProofTools:nonliftable-or-crossing}
\level{coarse}
Consider finite witnesses for the same $f,x,y,n$ on pages $r_{\rm early}\leq r_{\rm late}$. If the late witness is essential and the early one is inessential, then either a shorter early extension retains the original target and has a source in the earlier cycle set but outside the later cycle set, or the same late normalized extension spectral sequence has a nonzero higher-source crossing in the specified range. The first branch uses the actual correction label with $b=0$ and all length and cycle bounds.
% Blueprint note: This disjunction is distinct from the set-difference equivalence printed in Proposition~\ref{paper:6.20}.
\lean{KIP126.ProofTools.Paper.finiteExtension_nonliftableShort_or_lateCrossing}
\end{theorem}

\section{Generalized Leibniz at an explicit lambda level}

\begin{definition}[The normalized finite source at a specified level]
\label{sem:ProofTools:level-source-label}
\level{coarse}
Following the arbitrary-level quotient windows preceding Definition 5.4 of \cite{LWX126}, fix a quotient exponent $q$, a normalization shift $e$, a level $k\in\mathbb N$, and $x\in E_2^{s,t}(X)$. The normalized finite source label is the image of the native label $\lambda^k x$ under the actual $(0,e)$ shift and the inverse quotient-shift comparison, in tridegree
\[(s,t,t+e-k)\quad\text{of}\quad (\Sigma^{0,e}\nu X)/\lambda^q.\]
The chosen comparison respects both quotient inclusion and absolute boundary. This specifies an actual source label without assuming a new calibration of the existing comparison at positive shift and positive level.
\lean{KIP126.ProofTools.Standard.levelFiniteSourceLabel}
\end{definition}

\begin{definition}[A finite extension retaining its lambda level]
\label{sem:ProofTools:level-finite-extension}
\level{coarse}
Let $f:X\to Y$, $e=e(f)$, $A=\Sigma^{0,e}\nu X$, $B=\nu Y$, and $F:A\to B$ be the normalized map. For $n,k\in\mathbb N$, put $q=n+k-1$. A level-$k$ finite extension of length $m$ from $x\in E_2^{s,t}(X)$ to $y\in E_2^{s+m,t+m}(Y)$ consists of the bounds
\[2\le n,\qquad e\le m\le n-2+e,\]
and actual infinite-page source and target classes, represented by the normalized label $\lambda^k x$ and the quotient projection of $\lambda^{k+m-e}y$, which satisfy
\[d_m^{F_q}(\lambda^k x)=\lambda^{k+m-e}y\]
at fixed stem and weight $(t-s,t+e-k)$. The respective native windows are
\[Z_{n-1}^{s,t}(X)/B_{k+1}^{s,t}(X),\qquad
Z_{n-1-m+e}^{s+m,t+m}(Y)/B_{k+m-e+1}^{s+m,t+m}(Y).\]
The extension need not be essential, and zero labels are allowed.
\lean{KIP126.ProofTools.Standard.LevelFiniteExtension}
\uses{sem:ProofTools:level-source-label}
\end{definition}

\begin{definition}[No crossing on the three actual level edges]
\label{sem:ProofTools:level-no-crossing}
\level{coarse}
For an actual filtered map $g$ and fixed stem/weight $p$, write $NC(g;p,b,h)$ if there is no essential extension $d_{v-u}^g(\xi)=\eta$ for any $b<u\le v\le b+h$ and any associated-graded classes $\xi,\eta$. Essentiality is taken modulo the complete target indeterminacy of that actual map. With $A,B,F,q$ as above and $p=(t-s,t+e-k)$, the three predicates are respectively
\[
NC(F_q;p,s,m),\qquad NC(\delta_q^A;p,s,r),\qquad
NC(\Sigma^{1,-q}F;p,s+r,l).
\]
They refer to the top, left and bottom maps in the quotient-boundary square. The bottom condition is tested on the propagated infinite extension, at its actual level $r-n$; it is not inferred from a zero-level no-crossing condition.
\lean{KIP126.ProofTools.Standard.LevelFiniteNoCrossing,KIP126.ProofTools.Standard.LevelBoundaryNoCrossing,KIP126.ProofTools.Standard.LevelBottomNoCrossing}
\end{definition}

\begin{theorem}[Generalized Leibniz with an arbitrary lambda level]
\label{sem:ProofTools:generalized-leibniz-at-level}
\level{coarse}
Adapted from the quotient-boundary square proof of Theorem 6.1 of \cite{LWX126}, with its crossing hypotheses imposed on the actual level edges. In the fixed model, let $f:X\to Y$, $e=e(f)$, and
\[2\le n\le r,\qquad e\le m\le n-2+e,\qquad l\ge e.\]
Let $k\in\mathbb N$ be arbitrary, $q=n+k-1$, and $p=(t-s,t+e-k)$. Take
\[
x\in E_2^{s,t}(X),\quad y\in E_2^{s+m,t+m}(Y),\quad
x_\infty\in Z_\infty^{s+r,t+r-1}(X),\quad
y_\infty\in Z_\infty^{s+r+l,t+r+l-1}(Y).
\]
Suppose $d_r(x)=x_\infty$, there is the level-$k$ finite extension
\[d_m^{F_q}(\lambda^k x)=\lambda^{k+m-e}y,\]
and the existing zero-level infinite extension $d_l^F(x_\infty)=\lambda^{l-e}y_\infty$. Suppose also
\[
\bigl(NC(\delta_q^A;p,s,r)\lor NC(F_q;p,s,m)\bigr)
\land NC(\Sigma^{1,-q}F;p,s+r,l).
\]
Then the classical Adams differential is
\[d_{r+l-m}(y)=y_\infty.\]
There is no upper bound on $k$. The left edge and the propagated bottom edge are constructed in the proof; they are not extra caller premises.
\lean{KIP126.ProofTools.Standard.generalizedLeibnizAtLevel}
\uses{sem:ProofTools:level-finite-extension,sem:ProofTools:level-no-crossing}
% Blueprint note: This separate public theorem has a complete local proof relative to the registered upstream construction and proof obligations. It does not discharge either hole in the frozen original public GL, and does not identify the old crossing predicates with these actual conditions.
\end{theorem}

\begin{theorem}[The original finite extension supplies level zero]
\label{sem:ProofTools:finite-extension-at-level-zero}
\level{coarse}
Every original finite extension $d_m^{f,E_n}(x)=y$ supplies the level-zero finite extension in quotient $n-1$. The actual map equation and its prescribed source and target are retained through the proved representative comparisons. No additional cycle or crossing premise is required. This result does not compare the two sets of no-crossing hypotheses.
\lean{KIP126.ProofTools.Standard.finiteExtensionAtLevelZero}
\uses{sem:ProofTools:level-finite-extension}
\end{theorem}

\section{Conditional division by lambda}

\begin{definition}[The additional target-kernel crossing condition]
\label{sem:ProofTools:finite-lambda-kernel-no-crossing}
\level{coarse}
Let $f:X\to Y$, $e=e(f)$, $A=\Sigma^{0,e}\nu X$, $B=\nu Y$, and $F:A\to B$ be the normalized map. Fix $0<i<q$, $p=(t-s,t+e)$ and a target filtration $N$. In the actual lambda square put
\[
A'=\Sigma^{0,-i}(A/\lambda^{q-i}),\quad
B'=\Sigma^{0,-i}(B/\lambda^{q-i}),\quad
F'=\Sigma^{0,-i}F_{q-i}.
\]
Let $L_A:A'\to A/\lambda^q$ and $L_B:B'\to B/\lambda^q$ be the actual finite lambda maps. The extra crossing condition states that, for every
\[s+i\le u\le v<N,\]
an essential extension $d_{v-u}^{F'}(\xi)=\eta$ at $p$ has target satisfying
\[\operatorname{gr}_v(L_B)(\eta)\ne0.\]
The initial source row $u=s+i$ is included, while the terminal target row $v=N$ is excluded. This rules out essential smaller-layer targets killed by lambda; it neither requires injectivity on the entire graded target nor follows merely from absence of positive-length extensions of the lambda map itself.
\lean{KIP126.ProofTools.Standard.FiniteLambdaKernelNoCrossing}
\end{definition}

\begin{theorem}[Finite lambda division preserving the actual source and image]
\label{sem:ProofTools:finite-lambda-division}
\level{coarse}
Use the objects, maps, stem and weight of the preceding definition. Assume $s+i\le N$ and the additional target-kernel crossing condition. If an actual representative
\[a\in F^{s+i}\pi_{t-s,t+e}(A/\lambda^q)\]
has $F_q(a)\in F^N\pi_{t-s,t+e}(B/\lambda^q)$, there exists
\[b\in F^{s+i}\pi_{t-s,t+e}(A')\]
such that
\[L_A(b)=a,\qquad L_B(F'(b))=F_q(a),\qquad F'(b)\in F^N\pi_{t-s,t+e}(B').\]
Thus division preserves the specified actual source and image and reflects the requested target filtration. The smaller-layer source class is the class of the constructed $b$; no independently selected smaller source is prescribed, and kernel ambiguity at the final target grade is allowed.
\lean{KIP126.ProofTools.Standard.finiteLambdaDivision}
\uses{sem:ProofTools:finite-lambda-kernel-no-crossing}
% Blueprint note: The proof uses a finite correction over source and target filtrations, followed by the existing actual finite lambda source lift. It adds no separation or completeness hypothesis, and retains registered upstream obligations. The corresponding private essential-extension comparison maps both constructed smaller classes back to the specified larger classes.
\end{theorem}
